\documentclass[11pt]{article}

\usepackage[margin=0.92in]{geometry}
\usepackage{amsmath,amssymb,amsthm,mathtools,booktabs,enumitem,microtype}
\usepackage{xcolor}
\usepackage[colorlinks=true,linkcolor=blue!50!black,urlcolor=blue!50!black,citecolor=blue!50!black]{hyperref}

\theoremstyle{plain}
\newtheorem{theorem}{Theorem}
\newtheorem{lemma}[theorem]{Lemma}
\newtheorem{proposition}[theorem]{Proposition}
\newtheorem{corollary}[theorem]{Corollary}

\theoremstyle{definition}
\newtheorem{definition}[theorem]{Definition}

\theoremstyle{remark}
\newtheorem{remark}{Remark}

\DeclareMathOperator{\tr}{tr}
\DeclareMathOperator{\rank}{rank}
\DeclareMathOperator{\spn}{span}

\title{\bfseries Refutation of the Compact-Group Sensing Conjecture}
\author{}
\date{\vspace{-1.6em}%
\small
\begin{tabular}{@{}c@{}}
Target paper: J.~Tachella, D.~Chen, and M.~Davies,\\
\emph{Sampling Theorems for Unsupervised Learning in Linear Inverse Problems},\\
\href{https://arxiv.org/abs/2203.12513}{arXiv:2203.12513} (\emph{JMLR} 24(191):1--45, 2023), \S4.3, Conjecture~10\\[0.35em]
Companion verification: \texttt{solutions/unsupervised\_group\_sensing/explore.py}\\[0.35em]
November 2026
\end{tabular}}

\begin{document}

\maketitle

\begin{abstract}
We disprove Conjecture~10 of Tachella, Chen, and Davies (\href{https://arxiv.org/abs/2203.12513}{arXiv:2203.12513}; \emph{JMLR} 2023, \texttt{conj: all groups}), including its single-operator $\mathcal{G}$-invariant clause, under both the unbounded conic branch \textbf{A1(b)} and the bounded branch \textbf{A1(a)}.
A faithful diagonal sign action of $G = (\mathbb{Z}/2)^6$ (realized as $\mathbb{F}_2^3 \times \mathbb{F}_2^3$) on $\mathbb{R}^{28}$ decomposes into $28$ distinct one-dimensional characters ($c_j = s_j = 1$, $C = \max_j c_j/s_j = 1$) with exact nonidentity trace $\tr T_{(u,w)} = -4(-1)^{u\cdot w}$, so every nonidentity element has $16$ diagonal entries of one sign and $12$ of the other.
For the $G$-invariant set of $12$-sparse signals ($k = 12$), Conjecture~10 predicts unique generic identification for $m > 2k + C + 1 = 26$ (in particular $m = 27$), whereas the exact generic identification threshold is $m = 28$.
Generalizing to $\mathbb{F}_2^r \times \mathbb{F}_2^r$ ($r \ge 3$) produces an arbitrarily large gap of $2^{r-1} - 3$ measurements while $C = 1$ stays fixed.
We also determine exact generic identification thresholds for four comparison subspace models ($SO(2)$, $O(2)$, $SO(3)$, and $S_4$) via quotient Jacobian rank calculations.
\end{abstract}

\section{Problem Statement and Context}

Let $\mathcal{G}$ be a compact group acting on $\mathbb{R}^n$ via linear operators $\{T_g\}_{g \in \mathcal{G}}$, and let $\mathcal{X} \subset \mathbb{R}^n$ be a signal set fixed before choosing a linear measurement operator $A \in \mathbb{R}^{m \times n}$.
The observed data are the indexed measurement sets $\{A T_g \mathcal{X}\}_{g \in \mathcal{G}}$; when $\mathcal{X}$ is $\mathcal{G}$-invariant, every view equals $A\mathcal{X}$.
Unique identification of $\mathcal{X}$ means that no distinct admissible model $\mathcal{X}'$ satisfies $A T_g \mathcal{X}' = A T_g \mathcal{X}$ for all $g \in \mathcal{G}$.

\textbf{Assumption~A1} requires either:
\begin{enumerate}[label=\textup{\textbf{A1(\alph*)}},leftmargin=4.2em,itemsep=1pt]
\item $\mathcal{X} \subset \mathbb{R}^n$ is bounded with upper box-counting dimension $k$; or
\item $\mathcal{X} \subset \mathbb{R}^n$ is an unbounded cone whose intersection $\mathcal{X} \cap \mathbb{S}^{n-1}$ with the unit sphere has upper box-counting dimension $k - 1$.
\end{enumerate}

Over $\mathbb{C}$, the representation decomposes into $c_j$ copies of each occurring irreducible representation of complex dimension $s_j$, so $n = \sum_j c_j s_j$; set $C := \max_j c_j / s_j$.
Tachella, Chen, and Davies establish the compact-group necessary condition $m \ge C$ and the sufficient condition $m > 2k + C + 1$ for finite cyclic groups (and circle actions by restriction to finite cyclic subgroups preserving distinct weights).
Applying that theorem to a general compact group via a cyclic subgroup requires using the multiplicities of the \emph{restricted} representation; for independent generic operators, a separate theorem gives $m > k + n/|\mathcal{G}|$.
Conjecture~10 instead proposes replacing the cyclic restriction by the full-group irreducibles:

\begin{quote}
\textbf{Conjecture~10 (\texttt{conj: all groups}, \href{https://arxiv.org/abs/2203.12513}{arXiv:2203.12513} / \emph{JMLR} 2023, \S4.3).}
\emph{Let $\mathcal{G}$ be any compact group acting on $\mathbb{R}^n$ and assume the signal set $\mathcal{X}$ verifies Assumption~\textup{\textbf{A1}}. For almost every $A \in \mathbb{R}^{m \times n}$ with respect to Lebesgue measure, $\mathcal{X}$ can be uniquely identified from $\{A T_g \mathcal{X}\}_{g \in \mathcal{G}}$ if $m > 2k + \max_j c_j/s_j + 1$. In particular, for almost any $A \in \mathbb{R}^{m \times n}$, a $\mathcal{G}$-invariant signal set can be identified from $A\mathcal{X}$ if $m > 2k + \max_j c_j/s_j + 1$.}
\end{quote}

\section{Sign Representation with Small Minority Eigenspaces}

Let $G := \mathbb{F}_2^3 \times \mathbb{F}_2^3 \cong (\mathbb{Z}/2)^6$ under coordinatewise addition modulo $2$, and index the coordinates of $\mathbb{R}^{28}$ by the subset
\[
D := \{(a, b) \in \mathbb{F}_2^3 \times \mathbb{F}_2^3 : a \cdot b = 1\},
\]
where $a \cdot b \in \mathbb{F}_2$ denotes the standard dot product modulo $2$.
There are $2^3 - 1 = 7$ nonzero vectors $a \in \mathbb{F}_2^3$, and for each nonzero $a$ the affine hyperplane $\{b \in \mathbb{F}_2^3 : a \cdot b = 1\}$ has cardinality $4$, giving $n := |D| = 28$.
For each $g = (u, w) \in G$, define the diagonal orthogonal operator $T_g \in O(28)$ with entries
\[
(T_{(u,w)})_{(a,b),(a,b)} := (-1)^{u \cdot a + w \cdot b}
\qquad \bigl((a, b) \in D\bigr).
\]
Because the $28$ pairs $(a, b) \in D$ are distinct elements of $\mathbb{F}_2^6$, these are $28$ pairwise distinct real sign characters of $G$, remaining distinct one-dimensional irreducible representations after complexification.
Hence $c_j = s_j = 1$ for every occurring irreducible, and $C = \max_j c_j/s_j = 1$.

\begin{lemma}[Exact trace formula and faithfulness]\label{lem:trace}
For every nonidentity element $g = (u, w) \in G \setminus \{(0, 0)\}$,
\[
\tr T_{(u,w)} = -4(-1)^{u \cdot w}.
\]
In particular, every nonidentity $T_g$ has $16$ diagonal entries of one sign and $12$ of the opposite sign, so the representation is faithful.
\end{lemma}

\begin{proof}
Since $g = (u, w) \neq (0, 0)$, the character sum $\sum_{(a,b) \in G} (-1)^{u \cdot a + w \cdot b}$ vanishes.
Substituting the indicator $\mathbf{1}_{a \cdot b = 1} = \frac{1}{2}\bigl(1 - (-1)^{a \cdot b}\bigr)$ yields
\[
\tr T_{(u,w)}
= \sum_{(a,b) \in G} (-1)^{u \cdot a + w \cdot b}\,\frac{1 - (-1)^{a \cdot b}}{2}
= -\frac{1}{2} \sum_{a \in \mathbb{F}_2^3} (-1)^{u \cdot a} \sum_{b \in \mathbb{F}_2^3} (-1)^{(a + w) \cdot b}.
\]
The inner sum over $b \in \mathbb{F}_2^3$ equals $8$ if $a = w$ and $0$ otherwise, giving $\tr T_{(u,w)} = -4(-1)^{u \cdot w}$.
Thus for $g \neq 0$, the diagonal of $T_g$ consists of $16$ entries of one sign and $12$ of the other (specifically, $28$ elements have trace $+4$, $35$ have trace $-4$, and $T_0 = I$ has trace $28$).
Since $T_g \neq I$ for all $g \neq 0$, the action is faithful.
\end{proof}

For each $g \in G$, let $\epsilon_g \in \{+1, -1\}$ be the majority sign of $T_g$ (with $\epsilon_0 := +1$).
Then $T_g - \epsilon_g I$ vanishes on all majority-sign coordinates, so its image is supported on at most $12$ coordinates.

\section{Refutation of Conjecture 10 and Arbitrarily Large Gap}

\begin{theorem}[Sharp counterexample in $\mathbb{R}^{28}$]\label{thm:main}
Under the faithful $G = (\mathbb{Z}/2)^6$ action on $\mathbb{R}^{28}$ ($C = 1$), both the unbounded $G$-invariant cone $\mathcal{X} := \{x \in \mathbb{R}^{28} : \|x\|_0 \le 12\}$ (\textup{\textbf{A1(b)}}, $k = 12$) and its bounded truncation $\mathcal{X}_B := \mathcal{X} \cap B_1(0)$ (\textup{\textbf{A1(a)}}, $k = 12$) satisfy $2k + C + 1 = 26$, yet fail to be uniquely identified from $A\mathcal{X}$ (equivalently from $\{A T_g \mathcal{X}\}_{g \in G}$) for almost every $A \in \mathbb{R}^{m \times 28}$ whenever $m \le 27$.
For $m \ge 28$, almost every $A$ uniquely identifies both models.
Hence the exact generic identification threshold is $m = 28$, disproving Conjecture~10 at $m = 27$.
\end{theorem}

\begin{proof}
We first treat the unbounded conic branch \textbf{A1(b)}.
Since each $T_g$ is diagonal with $\pm 1$ entries, $\mathcal{X} = \{x \in \mathbb{R}^{28} : \|x\|_0 \le 12\}$ is a closed $G$-invariant cone.
Its spherical section $\mathcal{X} \cap \mathbb{S}^{27}$ is the finite union of $\binom{28}{12}$ coordinate $11$-spheres, which has upper box-counting dimension $11 = k - 1$, so $k = 12$.

Fix $m \le 27$.
For each coordinate $i \in \{1, \dots, 28\}$, deleting column $i$ from $A \in \mathbb{R}^{m \times 28}$ leaves a matrix of full rank $m$ outside the proper algebraic hypersurface where all $m \times m$ minors of the remaining $27$ columns vanish (which is proper because a rank-$m$ example exists).
On the complement of these $28$ measure-zero sets, column $i$ is a linear combination of the remaining columns for every $i$, so each coordinate functional $v \mapsto v_i$ is nonzero on $\ker A$.
Because a nonzero real vector space is not a finite union of proper hyperplanes, almost every $A \in \mathbb{R}^{m \times 28}$ admits a kernel vector $v \in \ker A$ with all $28$ coordinates nonzero.

For any such full-support $v \in \ker A$, define for each $g \in G$ the witness
\[
y_g := (T_g - \epsilon_g I) v.
\]
By Lemma~\ref{lem:trace}, $\|y_g\|_0 \le 12$ ($12$ for $g \neq 0$ and $0$ for $g = 0$), so $y_g \in \mathcal{X}$.
Because $Av = 0$,
\[
A T_g v = A(T_g - \epsilon_g I) v + \epsilon_g A v = A y_g \in A\mathcal{X}.
\]
Consider the competitor model
\[
\mathcal{X}'_A := \mathcal{X} \cup \bigcup_{g \in G} \spn_{\mathbb{R}}\{T_g v\}.
\]
Since $v$ has $28 > 12$ nonzero entries, $v \in \mathcal{X}'_A \setminus \mathcal{X}$, so $\mathcal{X}'_A \neq \mathcal{X}$.
The $64$ adjoined lines are permuted by $G$, so $\mathcal{X}'_A$ is a $G$-invariant cone whose spherical section adds at most $128$ isolated points to $\mathcal{X} \cap \mathbb{S}^{27}$, preserving upper box-counting dimension $11$ ($k = 12$).
For every $t \in \mathbb{R}$ and $g \in G$, $A(t T_g v) = A(t y_g) \in A\mathcal{X}$, while $\mathcal{X} \subset \mathcal{X}'_A$ gives the reverse inclusion; hence $A T_h \mathcal{X}'_A = A\mathcal{X}'_A = A\mathcal{X} = A T_h \mathcal{X}$ for all $h \in G$.
(Testing uniqueness of the fixed model $\mathcal{X}$ at a given $A$ allows the indistinguishable competitor $\mathcal{X}'_A$ to depend on $A$.)

For the bounded branch \textbf{A1(a)}, set $\mathcal{X}_B := \mathcal{X} \cap B_1(0)$, normalize $v \in \ker A$ so that $\|v\|_2 \le 1/2$, and define $\mathcal{X}'_{B,A} := \mathcal{X}_B \cup \{T_g v : g \in G\}$.
Each witness satisfies $\|y_g\|_2 = \|(T_g - \epsilon_g I)v\|_2 \le 2\|v\|_2 \le 1$, hence $y_g \in \mathcal{X}_B$ and $A T_g v = A y_g \in A\mathcal{X}_B$.
Adjoining finitely many points preserves upper box-counting dimension $12$, giving distinct bounded $G$-invariant \textbf{A1(a)} models with identical measurements.

Conversely, for every $m \ge 28$, almost every $A \in \mathbb{R}^{m \times 28}$ has rank $28$ and thus a left inverse, which recovers any subset of $\mathbb{R}^{28}$ uniquely from its image under $A$.
\end{proof}

\begin{corollary}[Arbitrarily large gap at $C = 1$]\label{cor:gap}
For every integer $r \ge 3$, let $G_r := \mathbb{F}_2^r \times \mathbb{F}_2^r$ act diagonally on $\mathbb{R}^{n_r}$ indexed by $D_r := \{(a, b) \in G_r : a \cdot b = 1\}$, where $n_r = (2^r - 1)2^{r-1}$.
Then $C = 1$, $\tr T_{(u,w)} = -2^{r-1}(-1)^{u\cdot w}$ for all $(u, w) \neq (0, 0)$, and for sparsity $k_r := (n_r - 2^{r-1})/2$ the exact generic identification threshold is $n_r$, exceeding the first integer $2k_r + C + 2 = n_r - 2^{r-1} + 3$ allowed by Conjecture~10 by $2^{r-1} - 3$.
\end{corollary}

\begin{proof}
There are $2^r - 1$ nonzero vectors $a \in \mathbb{F}_2^r$, each with $2^{r-1}$ partners $b$ satisfying $a \cdot b = 1$, so $n_r = (2^r - 1)2^{r-1}$.
The character-sum computation of Lemma~\ref{lem:trace} yields $\tr T_{(u,w)} = -\frac{1}{2}\cdot 2^r (-1)^{u\cdot w} = -2^{r-1}(-1)^{u\cdot w}$ for $(u, w) \neq (0, 0)$.
Hence every nonidentity element has minority-sign count $(n_r - 2^{r-1})/2 = k_r$, and the argument of Theorem~\ref{thm:main} shows that generic identification fails for all $m < n_r$ and holds for $m \ge n_r$.
The first integer strictly greater than $2k_r + C + 1 = 2k_r + 2$ is $2k_r + 3 = n_r - 2^{r-1} + 3$, giving a gap of $2^{r-1} - 3$.
\end{proof}

\section{Comparison Subspace Models}

To contrast how the real linear span of the group representation governs identification beyond the normalized multiplicity $C$, let $W$ be a $d$-dimensional irreducible real representation of a compact group $\mathcal{G}$, consider the diagonal action on $W^6$ (ambient dimension $6d$), and fix the $\mathcal{G}$-invariant conic model $\mathcal{X} := W \oplus 0^5$ of dimension $k = d$.

\begin{proposition}[Quotient Jacobian thresholds on $W^6$]\label{prop:subspace}
Among $\mathcal{G}$-invariant \textup{\textbf{A1}} models of dimension $d$, the exact generic identification threshold for $\mathcal{X} = W \oplus 0^5$ is $m = 7$ for $SO(2)$ on $\mathbb{R}^2$ ($d = 2$), and $m = 5$ for $O(2)$ on $\mathbb{R}^2$ ($d = 2$), $SO(3)$ on $\mathbb{R}^3$ ($d = 3$), and $S_4$ on its standard sum-zero representation ($d = 3$).
\end{proposition}

\begin{proof}
If $m < d$, generically $A\mathcal{X} = \mathbb{R}^m$, so adjoining any hidden copy of $W$ preserves $A\mathcal{X}$ and dimension $d$.
Suppose $m \ge d$. Generically $A|_{\mathcal{X}}$ is injective; quotienting $\mathbb{R}^m$ by $A\mathcal{X}$ yields $q := m - d$ residual linear rows $B \in \mathbb{R}^{q \times 5d}$ acting on the five hidden copies $W^5$.
Conditional on any full-rank visible block, the map from the hidden columns of $A$ to $B$ is a linear surjection in fixed quotient coordinates, so generic rank tests on $B$ apply directly.

A vector $(x_0, x_{\mathrm{hid}}) \in W \oplus W^5$ satisfies $A T_g(x_0, x_{\mathrm{hid}}) \in A\mathcal{X}$ for all $g \in \mathcal{G}$ if and only if $x_{\mathrm{hid}}$ belongs to the common kernel $K := \bigcap_{g \in \mathcal{G}} \ker\bigl(B(I_5 \otimes T_g|_W)\bigr)$.
Thus $\mathcal{X} \oplus K$ is the largest $\mathcal{G}$-invariant set consistent with $A\mathcal{X}$.
If $K = \{0\}$, every competitor lies in $\mathcal{X}$ and injectivity of $A|_{\mathcal{X}}$ forces equality.
If $K \neq \{0\}$, $K$ is a nonzero $\mathcal{G}$-invariant subspace of $W^5$, so by complete reducibility it contains an irreducible copy $W' \subset K$ of $W$; then $\mathcal{X} \cup W'$ is a distinct $\mathcal{G}$-invariant cone of dimension $d$ with image $A\mathcal{X}$.
Hence generic identification holds if and only if the stacked quotient Jacobian $\{B(I_5 \otimes T_g|_W)\}_{g \in \mathcal{G}}$ has full rank $5d$.

For $SO(2)$ on $\mathbb{R}^2$, $\spn_{\mathbb{R}}\{T_g|_W\} = \spn_{\mathbb{R}}\{I_2, J\}$, where $J$ is the quarter-turn.
On $\mathbb{C}^5 \simeq W^5$, each real probe and its $J$-rotate span a complex line in $(\mathbb{C}^5)^*$, giving generic real rank $2\min(q, 5)$; full rank $10$ first occurs at $q = 5$ ($m = 7$).
For $O(2)$, $SO(3)$, and $S_4$ (sum-zero), the group operators span all of $M_d(\mathbb{R})$ (certified on finite test subsets below).
Reshaping each row of $B$ as a $d \times 5$ matrix, left multiplication by matrix units shows that its $M_d(\mathbb{R})$-translates span all $d \times 5$ matrices whose rows lie in that probe's row space.
The combined row spaces of $q$ generic probes in $\mathbb{R}^5$ have dimension $\min(dq, 5)$, giving stacked rank $d\min(dq, 5)$.
Full rank $5d$ first occurs at $q = \lceil 5/d \rceil$: $q = 3$ ($m = 5$) for $O(2)$ and $q = 2$ ($m = 5$) for $SO(3)$ and $S_4$.
\end{proof}

\section{Exact Computational Verification}

The companion script \texttt{explore.py} in \texttt{solutions/unsupervised\_group\_sensing/} executes exact integer and SymPy arithmetic via \texttt{.venv/bin/python -B}:
\begin{itemize}[leftmargin=2em,itemsep=1pt,topsep=3pt]
\item Verifies $M^\top M = 64 I_{28}$ for the $64 \times 28$ character table of $G = (\mathbb{Z}/2)^6$ on $D$, and enumerates all $64$ traces: $\{28: 1,\; 4: 28,\; -4: 35\}$.
\item Constructs an explicit rank-$27$ integer sensor $A \in \mathbb{Z}^{27 \times 28}$ with full-support kernel vector $v = (1, 2, \dots, 27, 1)^\top$, verifies all $64$ witnesses $y_g = (T_g - \epsilon_g I)v$ (sparsity $12$ for $g \neq 0$, $0$ for $g = 0$) and $A T_g v = A y_g$, confirms residual Jacobian blocks $\epsilon_g A$ (rank $27$), and verifies rank $28$ upon adding one independent sensor row.
\item Certifies the full action-algebra span dimensions ($2$ for $SO(2)$, $d^2$ for $O(2)$, $SO(3)$, and $S_4$) using exact finite test sets (signed permutations for $O(2)$, determinant-$1$ signed permutations for $SO(3)$, and all $24$ permutations on the sum-zero subspace for $S_4$) and verifies the quotient Jacobian ranks in Table~\ref{tab:subspace}.
\end{itemize}

\begin{table}[ht]
\centering
\small
\begin{tabular}{@{}lrrr@{}}
\toprule
Group and representation & Ambient dim. & Quotient Jacobian ranks ($q = 0, 1, \dots$) & First identifying $m$ \\
\midrule
$SO(2)$, plane ($d=2$)     & 12 & 0, 2, 4, 6, 8, 10, 10 & 7 \\
$O(2)$, plane ($d=2$)      & 12 & 0, 4, 8, 10, 10       & 5 \\
$SO(3)$, standard ($d=3$)  & 18 & 0, 9, 15, 15          & 5 \\
$S_4$, sum-zero ($d=3$)    & 18 & 0, 9, 15, 15          & 5 \\
\bottomrule
\end{tabular}
\caption{Exact quotient Jacobian ranks and generic identification thresholds on $W^6$ for $\mathcal{X} = W \oplus 0^5$, certified with exact integer arithmetic in \texttt{explore.py}.}
\label{tab:subspace}
\end{table}

\end{document}
